(1) With $(i\sigma)^{2} = -\sigma^{2} = -1$ the exponential series splits into cosine and sine, giving the displayed form; it is unitary because $i\sigma$ is skew-adjoint ($\sigma$ is self-adjoint, being $2N - 1$); $R_{2\pi i} = \cos\pi + i\sigma\sin\pi = -1$ and $R_{4\pi i} = +1$. (2) $\sigma = \operatorname{diag}(-1, 1)$ on the module (Theorem \claimref{singularity.descent-parity-average}'s proof), so $e^{i\theta\sigma/2} = \operatorname{diag}(e^{-i\theta/2}, e^{i\theta/2})$, which is $-1$ at $\theta = 2\pi$; conjugation by a scalar is trivial. (3) The even part of $\Cl(1,3)$ is spanned by $1$, the six frame bivectors and $I$, all with real coefficients, and $i\sigma_r$ is none of these: the bivectors of $\Cl(1,3)$ with square $-1$ are the spatial planes. $I\sigma_r = e_2e_3$ is one line of blade multiplication, and a rotor of the plane $e_2e_3$ fixes every vector orthogonal to that plane, $e_0$ and $e_1$ included, hence fixes $\sigma_r = e_0e_1$; the script verifies $I\sigma_r = e_2e_3$, the action of its rotor on all four frame vectors, and the unitary rotor's algebra on the module, blade-exactly.
